%%
%% This is file `sample-sigconf-authordraft.tex',
%% generated with the docstrip utility.
%%
%% The original source files were:
%%
%% samples.dtx  (with options: `all,proceedings,bibtex,authordraft')
%% 
%% IMPORTANT NOTICE:
%% 
%% For the copyright see the source file.
%% 
%% Any modified versions of this file must be renamed
%% with new filenames distinct from sample-sigconf-authordraft.tex.
%% 
%% For distribution of the original source see the terms
%% for copying and modification in the file samples.dtx.
%% 
%% This generated file may be distributed as long as the
%% original source files, as listed above, are part of the
%% same distribution. (The sources need not necessarily be
%% in the same archive or directory.)
%%
%%
%% Commands for TeXCount
%TC:macro \cite [option:text,text]
%TC:macro \citep [option:text,text]
%TC:macro \citet [option:text,text]
%TC:envir table 0 1
%TC:envir table* 0 1
%TC:envir tabular [ignore] word
%TC:envir displaymath 0 word
%TC:envir math 0 word
%TC:envir comment 0 0
%%
%% The first command in your LaTeX source must be the \documentclass
%% command.
%%
%% For submission and review of your manuscript please change the
%% command to \documentclass[manuscript, screen, review]{acmart}.
%%
%% When submitting camera ready or to TAPS, please change the command
%% to \documentclass[sigconf]{acmart} or whichever template is required
%% for your publication.
%%
%%
\documentclass[manuscript,screen,review,anonymous]{acmart}

% ---------- Common, venue-agnostic preamble ----------
\usepackage[utf8]{inputenc}
\usepackage[english]{babel}

% Math / graphics
%\usepackage{amsmath,amssymb}
\usepackage{graphicx}
\DeclareGraphicsExtensions{.pdf,.png,.jpg}

% Tables / lists
\usepackage{booktabs}
\usepackage{multirow}
\usepackage{enumitem}
\usepackage{makecell}
\usepackage{relsize}

% Colors & drawings & boxes
\usepackage{xcolor}
\usepackage[most]{tcolorbox}
\usepackage[normalem]{ulem}
\usepackage{colortbl}

% Listings (portable)
\definecolor{darkgreen}{RGB}{0,100,0}
\definecolor{codegreen}{rgb}{0,0.5,0}
\definecolor{codepurple}{rgb}{0.58,0,0.82}
\definecolor{codegray}{rgb}{0.5,0.5,0.5}
\usepackage{listings}
\lstdefinestyle{mystyle}{
  commentstyle=\color{codegreen},
  keywordstyle=\bfseries,
  stringstyle=\color{codepurple},
  basicstyle=\ttfamily\scriptsize,
  breaklines=true,
  captionpos=b,
  keepspaces=true,
  tabsize=2
}
\lstset{style=mystyle}

% Cross-refs (safe as long as we don't load 'caption' ourselves)
%\usepackage[capitalize,nameinlink]{cleveref}

% Algorithms (you can drop this if unused)
\usepackage[noend,ruled,linesnumbered]{algorithm2e}
%\usepackage[ruled,vlined,linesnumbered]{algorithm2e}

% Hyphenation
\hyphenation{op-tical net-works semi-conduc-tor}

\usepackage{tikz}
\usetikzlibrary{arrows.meta,positioning,shapes,calc} % needed by the figure

% ---------- Common macros ----------
\newcommand{\find}[1]{%
\begin{tcolorbox}[tile,size=fbox,boxsep=2mm,boxrule=0pt,top=0pt,bottom=0pt,
borderline={0.6mm}{0pt}{black!66!white},colback=black!5!white]
\em #1
\end{tcolorbox}
}

\newcommand{\question}[1]{%
\begin{tcolorbox}[tile,size=fbox,boxsep=1.5mm,boxrule=0pt,top=0pt,bottom=0pt,
borderline west={1mm}{0pt}{blue!5!white},colback=blue!5!white]
\em #1
\end{tcolorbox}
}

\newcommand{\modify}[1]{\textcolor{blue}{#1}}
\newcommand{\changemark}[1]{{\color{black}#1}}
\newcommand{\redout}[1]{\sout{#1}}

\newcommand{\Honggfuzz}{\href{https://github.com/google/honggfuzz}{Honggfuzz}}
\newcommand{\AFLpp}{\href{https://github.com/AFLplusplus/AFLplusplus}{AFL++}}

\newcommand{\firstAuthor}{$^{\pdftooltip{\textrm{\faUser}}{first author}}$}
\newcommand{\corrAuthor}{$^{\pdftooltip{\textrm{\faEnvelope}}{corresponding author}}$}

\newcommand{\denseitems}{\setlength{\itemsep}{1pt}\setlength{\parskip}{0pt}}

\newcommand{\tclcolor}{\cellcolor[HTML]{C0C0C0}}

\newboolean{COMMENTSON} 
\setboolean{COMMENTSON}{true}   
\ifthenelse{\boolean{COMMENTSON}}
{
\newcommand{\WL}[2]{\todo[color=red!40,linecolor=red,bordercolor=red]{Wen: #2 } #1}
\newcommand{\wen}[1]{\todo[inline,color=red!40,linecolor=red,bordercolor=red]{Wen: #1 }}
}

\definecolor{DarkOrange}{rgb}{0.8,0.3,0.0} 
\definecolor{DarkCyan}{rgb}{0.0, 0.55, 0.55}
\definecolor{codegreen}{rgb}{0,0.6,0}
\definecolor{codegray}{rgb}{0.5,0.5,0.5}
\definecolor{codepurple}{rgb}{0.58,0,0.82}
\definecolor{backcolour}{rgb}{0.95,0.95,0.92}


% ---- Paper identity ----
\newcommand{\papertitle}{Beyond Source: An Empirical Study of Python Bytecode Security Risks}
\newcommand{\papertitleshort}{Python Bytecode Security Risks}

\newcommand{\tech}{\mbox{\textsc{PyBC-Sec}}}
\newcommand{\tool}{\tech}

% ---- Keywords ----
\newcommand{\paperkeywords}{Python bytecode, package security, software supply chain, program analysis, runtime robustness}

% ---- Authors ----
\newcommand{\authorAname}{Wen Li}
\newcommand{\authorAaffil}{Utah State University}
\newcommand{\authorAemail}{awen.li@usu.edu}

\newcommand{\authorBname}{Coauthor Name}
\newcommand{\authorBaffil}{Somewhere University}
\newcommand{\authorBemail}{coauthor@example.com}

% appendix for journal
\newcommand{\appendixcontent}{%
  \appendix
}


%%
%% \BibTeX command to typeset BibTeX logo in the docs
\AtBeginDocument{%
  \providecommand\BibTeX{{%
    Bib\TeX}}}

% Remove rights/publication metadata for submission
%\setcopyright{acmlicensed}
%\copyrightyear{2018}
%\acmYear{2018}
%\acmDOI{XXXXXXX.XXXXXXX}
%\acmConference[Conference acronym 'XX]{Make sure to enter the correct
%  conference title from your rights confirmation email}{June 03--05,
%  2018}{Woodstock, NY}
%\acmISBN{978-1-4503-XXXX-X/2018/06}


%%
%% Submission ID.
%% Use this when submitting an article to a sponsored event. You'll
%% receive a unique submission ID from the organizers
%% of the event, and this ID should be used as the parameter to this command.
%%\acmSubmissionID{123-A56-BU3}

%%
%% end of the preamble, start of the body of the document source.
\begin{document}


\title{\papertitle}

%%
%% the authors and their affiliations.
\author{\authorAname}
\affiliation{%
  \institution{\authorAaffil}
  \city{Logan}
  \country{USA}
}
\email{\authorAemail}

%%
%% The abstract is a short summary of the work to be presented in the article.
%!TEX root = paper.tex

\begin{abstract}
Python security analysis is largely source-centric: package scanners, dependency analyzers, malware detectors, and program-analysis pipelines typically inspect Python source files, metadata, installation scripts, and dependency relations.  However, CPython executes compiled code objects, routinely materializes bytecode as \texttt{.pyc} files, and can load bytecode through source-less modules, cached imports, bundled applications, custom loaders, and marshalled code objects.  This creates a practical visibility gap: behavior may be present in executable bytecode even when corresponding source is missing, stale, mismatched, or difficult to recover.

This paper presents an empirical study of Python bytecode as a first-class security artifact.  We study how bytecode appears in real software distributions, whether it is visible from source, how well existing tools recover source-level representations, and how CPython behaves when processing unusual bytecode inputs.  Our methodology combines package-scale artifact extraction, source/bytecode correspondence analysis, decompilation and recompilation checks, runtime robustness testing across CPython versions, and source-reachability analysis for abnormal behaviors.  The study is designed to separate ecosystem visibility risks from interpreter robustness findings and to distinguish bytecode-only behavior from behavior that ordinary source-level tooling can observe.

Our results show that Python bytecode deserves explicit treatment in package security workflows.  Bytecode artifacts occur in multiple distribution forms, source correspondence is not guaranteed, decompilation-based recovery is incomplete, and bytecode-level testing reaches runtime behaviors that are not always reproducible from ordinary source.  These findings motivate bytecode-aware package scanning, explicit source/bytecode mismatch checks, more robust bytecode validation boundaries, and security tooling that treats compiled Python artifacts as analyzable inputs rather than disposable caches.

\end{abstract}


%%
%% The code below is generated by the tool at http://dl.acm.org/ccs.cfm.
%% Please copy and paste the code instead of the example below.
%%
\begin{CCSXML}
<ccs2012>
   <concept>
       <concept_id>10002978.10003022</concept_id>
       <concept_desc>Security and privacy~Software and application security</concept_desc>
       <concept_significance>500</concept_significance>
       </concept>
   <concept>
       <concept_id>10011007.10011074.10011099.10011102.10011103</concept_id>
       <concept_desc>Software and its engineering~Software testing and debugging</concept_desc>
       <concept_significance>500</concept_significance>
       </concept>
 </ccs2012>
\end{CCSXML}

\ccsdesc[500]{Software and its engineering~Software testing and debugging}
\ccsdesc[500]{Security and privacy~Software and application security}

%%
%% Keywords. The author(s) should pick words that accurately describe
%% the work being presented. Separate the keywords with commas.
\keywords{configuration testing, redundancy reduction, program exploration, test selection}


%%
%% This command processes the author and affiliation and title
%% information and builds the first part of the formatted document.
\maketitle


%% The next two lines define the bibliography style to be used, and
%% the bibliography file.
\bibliographystyle{ACM-Reference-Format}
\bibliography{reference}


%%
%% If your work has an appendix, this is the place to put it.
\appendix



\end{document}
\endinput
%%
%% End of file `sample-sigconf-authordraft.tex'.
